% !TeX root = ../main.tex

在第~\ref{sec:preprocessing}~节处理结果上，结合变量定义、训练集分布与已观察的风险关系，
确定表~\ref{tab:initial-binning} 的初始分箱方案。连续变量采用便于解释的业务边界，
离散计数保留低次数差异并合并稀疏尾部；缺失及未知特殊编码优先按既定规则路由。
各变量独立分箱，同一客户可以同时属于收入特殊档和某逾期字段的特殊档，保留全部记录。

\begin{table}[!htbp]
  \centering
  \caption{十个原始特征的初始分箱方案}
  \label{tab:initial-binning}
  \footnotesize
  \setlength{\tabcolsep}{3pt}
  \renewcommand{\arraystretch}{1.12}
  \begin{tabular}{>{\raggedright\arraybackslash}p{2.2cm}>{\raggedright\arraybackslash}p{8.0cm}>{\raggedright\arraybackslash}p{4.3cm}}
    \toprule
    特征 & 候选分箱 & 主要考虑 \\
    \midrule
    循环额度使用率 & $0$；$(0,0.1]$；$(0.1,0.3]$；$(0.3,0.6]$；$(0.6,1]$；$(1,+\infty)$ & 零使用率单列；保留超额度状态。 \\
    年龄 & $(0,30)$；$[30,40)$；$[40,50)$；$[50,60)$；$[60,70)$；$[70,+\infty)$ & 按年龄段划分，考察风险随年龄的变化。 \\
    30--59 天逾期次数 & $0$；$1$；$2$；$3$；普通值 $\geq4$；96/98 特殊档 & 区分首次与重复逾期，合并普通高次数尾部。 \\
    负债收入比 & $0$；$(0,0.4]$；$(0.4,0.6]$；$(0.6,1]$；$(1,10]$；$(10,100]$；$(100,+\infty)$ & 保留 1 的业务边界及非单调尾部。 \\
    月收入 & 缺失；$[0,10)$；$[10,4{,}000)$；$[4{,}000,6{,}000)$；$[6{,}000,10{,}000)$；$[10{,}000,+\infty)$ & 前两档为既定特殊状态，普通收入再分层。 \\
    开放信贷账户数 & $0$；$1$--$2$；$3$--$5$；$6$--$9$；$10$--$14$；$\geq15$ & 兼顾信贷经历不足与账户较多的不同含义。 \\
    90 天及以上逾期次数 & $0$；$1$；$2$；普通值 $\geq3$；96/98 特殊档 & 保留严重逾期次数的风险梯度。 \\
    房地产贷款或额度数 & $0$；$1$--$2$；$3$；$\geq4$ & 区分无此类贷款、少量持有及较多持有。 \\
    60--89 天逾期次数 & $0$；$1$；普通值 $\geq2$；96/98 特殊档 & 普通 $\geq3$ 次仅 292 人，先与 2 次合并。 \\
    家属人数 & $0$；$1$；$2$；$3$；$\geq4$ & 基于填补后数值分箱，另保留原缺失指示。 \\
    \bottomrule
  \end{tabular}
\end{table}

\paragraph{分箱约束与标记用途}
普通箱以训练样本数不少于 450（训练集的 0.5\%）、好坏样本各不少于 30 作为初步检查标准，
结合相邻风险差异与业务含义决定是否合并。
具有明确识别规则的特殊档（例如借款人的三类逾期 96/98 编码）可保留较小样本规模。
家属人数的零值箱包含原生零值及缺失后补零的记录，
同时将 \texttt{NumberOfDependents\_missing} 的 0/1 两档作为候选输入，保留缺失信息。

\paragraph{DebtRatio 尾部分档的依据}
字典将 DebtRatio 定义为每月债务偿付、赡养费和生活费用与月总收入之比。
因此，大于 1 表示这些支出超过当月收入，提示负担较重、需要进一步核查，不能仅据此判为错误。
表~\ref{tab:initial-binning} 中大于 1 的各档仍为普通数值档，与未知编码的特殊档含义不同。
图~\ref{fig:training-feature-risk} 显示，坏样本率由 $(0.2,0.4]$ 档的 5.65\%
上升至 $(0.4,0.6]$、$(0.6,1]$ 和 $(1,10]$ 档的 7.81\%、10.28\% 和 10.50\%，
随后在 $(10,100]$ 和大于 100 档降至 4.34\% 和 5.97\%。
两个尾部档分别有 2,744 和 14,589 人，不能将其下降简单归为极少数记录造成的波动。

交叉统计进一步发现，大于 100 的 14,589 人中，收入缺失、为零和大于零但低于 10 的人数
分别为 13,586、724 和 276，三组合计 14,586 人，占 99.98\%；
相应坏样本率分别为 6.14\%、4.01\% 和 2.90\%。
DebtRatio 大于 1 后，各区间的坏样本率差异明显，且未呈现单调变化；极高值又集中于收入缺失或极低的样本，
不能简单理解为比率越大、风险越高。因此，保留这些区间分别分箱，避免强行合并掩盖风险差异，
后续再结合收入状态和验证集表现决定是否调整。
